\clearpage
\gdef\otChapNum{0}\pagestyle{otFront}\phantomsection\label{epreuve-guide}
\pdfbookmark[0]{Préparer l'épreuve anticipée}{bm-epreuve}
\otPageTitre{Préparer l'épreuve}{Première générale · Spécialité mathématiques · 2026-2027}
\fontsize{11}{15}\selectfont
\editionEtape{Le cadre officiel}
La note de service du 11 septembre 2026 définit l'épreuve pour l'année 2026-2027,
au titre du baccalauréat session 2028 ; elle remplace celle du 10 juin 2025.
L'épreuve écrite dure \textbf{2 heures}, a pour coefficient 2 et se passe
\textbf{sans calculatrice}. Le candidat de spécialité compose sur le sujet propre à cet enseignement.
\par\smallskip
{\renewcommand{\arraystretch}{1.12}\begin{tabularx}{\linewidth}{@{}p{43mm}Y@{}}\toprule
Première partie · 6 points & Automatismes : liste de questions à \textbf{réponses courtes}.\\
Deuxième partie · 14 points & Deux ou trois exercices indépendants.\\
Maîtrise de la langue & 2 points sur 20 : orthographe, syntaxe, lexique, organisation du propos.\\
Oral de contrôle & 20 minutes de préparation, 20 minutes d'entretien ; au moins deux exercices
sur des parties différentes du programme, sans calculatrice.\\\bottomrule
\end{tabularx}}\par\vspace{2mm}
\href{https://www.education.gouv.fr/bo/2026/Special4/MENE2622640N}{\textcolor{otNavy}{Définition officielle : note de service du 11 septembre 2026, BO spécial n° 4.}}
\editionEtape{Les automatismes évalués}
Ils relèvent de la rubrique « Automatismes » du programme de Première : évolutions et
variations, calcul numérique et algébrique, fonctions et représentations, statistiques,
probabilités ; s'y ajoutent les automatismes de Seconde. Le chapitre 14 et les sujets A et B
les entraînent. Les sujets zéro publiés en 2025 proposaient un QCM : pour la session 2028,
le candidat écrit lui-même chaque réponse courte.
\editionEtape{Deux entraînements originaux}
Les sujets A et B, créés pour ce livre, ne sont pas des sujets officiels. Chacun comporte
12 automatismes, puis deux exercices indépendants. Leur barème, choix pédagogique du livre, donne 0,5 point
par automatisme et 7 points par exercice, dont 1 point de maîtrise de la langue : orthographe,
phrases construites, vocabulaire précis, étapes enchaînées jusqu'à la conclusion.
\editionEtape{Une gestion du temps à adapter}
Environ 20 minutes pour les automatismes, 40 minutes par exercice et 20 minutes pour relire.
Après la séance, classer chaque erreur (cours, calcul, lecture, rédaction) et refaire plus tard
l'exercice de chapitre correspondant.
\editionRetenir{Une valeur exacte, un signe justifié, un domaine précisé et une phrase
qui répond à la question rendent le raisonnement vérifiable.}
\clearpage

